\section{Comprobación del cofactor para todas las caras}
Fijamos una máscara y un conjunto de $v$ sitios. Se quitan los $v$ del exterior
también en las caras que desactivan alguna de sus exclusiones. Supongamos el exterior
RC conectado, su punto estacionario regular y $0<x_i<1/d_i$ para las marcas presentes.
Todos los objetos en este cálculo se refieren a la rama difusa. Cada fibra RC marcada
tiene margen exterior $t_b=1-\sum_{i\in b}x_i>0$; la vacancia de una diagonal activa
es $q_f=1-\sum_{i\in f}x_i>0$. Para una fibra RC escribimos también $q_b=t_b$.
$U_R,U_D$ son las matrices de incidencias RC y diagonales de los sitios marcados.
El grado $d_i$ incluye ambas incidencias RC y sus diagonales locales activas.
Se pueden omitir diagonales vacías; un factor diagonal unitario tiene capacidad uno.
\paragraph{Identidad física.}
Sea $S_{\rm ext}$ el operador de líneas de torres sólo con las casillas exteriores.
Tomamos en las coordenadas RC el complemento ortogonal de $(\mathbf1,-\mathbf1)$.
Cada columna de $U_R$ es ortogonal a ese vector. En estas coordenadas, escribiendo
$U=(U_R^t,U_D^t)^t$, el operador físico y los factores de sitio son
\[
\begin{gathered} S=S_0+UTU^t,\quad S_0=\operatorname{diag}(S_{\rm ext}|_{\rm RC^\perp},
                    \operatorname{diag}(q_f:f\ \mathrm{diagonal})),\\
 T_{ii}=\frac{x_i(1-x_i)}{1-d_i x_i},\quad b_i=\frac{1-d_i x_i}{(1-x_i)^{d_i}}.\end{gathered}
\]
Aquí $\mathrm{RC}^\perp$ denota exclusivamente el complemento del modo nulo RC.
Como $S_0$ es invertible, el lema del determinante da, con $\prod_D$ sólo en diagonales,
\[
 \frac{\operatorname{pdet}S}{\operatorname{pdet}S_{\rm ext}\prod_D q_f}
   =\det T\det\!\left(T^{-1}+U_R^t S_{\rm ext}^+U_R+
                                    U_D^t\operatorname{diag}(1/q_f)U_D\right).
\]
Por la respuesta de los multiplicadores de margen del exterior, la Hessiana de
$R=H_{\rm ext}+\sum_b t_b\log t_b$ es
$R''=U_R^t\operatorname{diag}(1/t_b)U_R-U_R^t S_{\rm ext}^+U_R$.
En efecto, al variar las masas se resuelve la ecuación de márgenes en el mismo
complemento; el máximo exterior aporta la forma negativa de la inversa, y el
segundo término de $R$ aporta la forma positiva indicada. Cada factor categórico
local aporta $\operatorname{diag}(1/x_i)+{\bf1}{\bf1}^t/q_f$ a su Hessiana negativa.
Restando los $d_i-1$ términos de Bernoulli resulta
\[
 -H_B''=T^{-1}+U_R^t\operatorname{diag}(1/t_b)U_R
                         +U_D^t\operatorname{diag}(1/q_f)U_D.
\]
El último determinante físico es pues $\det(-H_B''-R'')$. Como
$b_i T_{ii}=x_i(1-x_i)^{1-d_i}$, al introducir $V_0=\operatorname{diag}(x_i(1-x_i))$
el factor local completo toma la forma
\[
 D_{\rm loc}=\det[(-H_B''-R'')V_0]\,
                    \prod_f\frac{q_f}{\prod_{i\in f}(1-x_i)} .
\]
La cantidad $\prod_f$ incluye RC. Su cociente con el producto que viene de $S$
es precisamente $\prod_b t_b$, cancelado en el logaritmo por
$+\frac12\sum_b\log t_b$ en $A_*$. El covolumen es el mismo porque el modo RC de
norma $\sqrt{2j}$ sigue siendo el único modo nulo. Las multiplicidades de una
fibra con varios sitios ya figuran en $U_R$ y en el mismo $t_b$.
\paragraph{Campos comunes y supresión física.}
El término afín extraído de $R$ sólo depende del exterior común. Al apagar una marca
poniendo $\zeta_i=0$, la rama pequeña da $x_i=0$. La columna $i$ de
$(-H_B''-R'')V_0$ tiende a $e_i$ y la submatriz complementaria es la de la cara.
Los cocientes escalares se restringen de igual modo si difieren las incidencias de
ese sitio. La partición exacta también ignora los monomios que contienen ese sitio.
Esto identifica la proyección en las dos evaluaciones, incluso en el normalizador.
Al contrario, para evaluar \emph{un átomo ocupado} se suprime una fila y una columna.
Se usa el cofactor exterior ya regularizado. Con $m=j-h$ supervivientes por lado,
en la notación de \emph{Media matriz}, parte 05, los grados escritos como $t_b d_b$
tienen $d_b\to1$ en índices apagados. En el camino $x=\lambda {\bf1}_X$,
la conductancia exterior es $\tau_2(p_{rc})$ con $p_{rc}=t_rt_c w_{rc}/j$,
donde $w$ permanece acotado hasta $\lambda=1$ por el balanceo en medidas $t\mu$.
La raíz para $w$ es regular y su jacobiano de margen conserva la brecha para
$h/j$ pequeño (parte 02, con las copias de índice de masa cero). Dividir por
$\sqrt{t_rd_rt_cd_c}$ deja $O(\sqrt{t_rt_c}/j)$ incluso si ambos índices se apagan.
Así las filas y columnas apagadas desaparecen también del núcleo normalizado.
La misma conductancia total vale $j\bar d_j=m\bar d_m$. El término
$\frac12\log(\bar d/2)$ cambia por $\frac12\log(m/j)$ respecto a la fórmula intrínseca
del menor si se excluye el covolumen en ambas fórmulas. Al incluirlo se obtiene
$\frac12\log(4\pi j)+\frac12\log(m/j)=\frac12\log(4\pi m)$.
En este límite se siguen las medidas y grados normalizados, sin tomar el
pseudodeterminante de una matriz con filas cero y el antiguo número de modos.
Se presupone $m>0$ y exterior superviviente conectado, asegurado por la brecha
en la parte 02. Sólo el prefactor exterior se evalúa en el átomo duro.
\paragraph{Alcance.}
El álgebra no tiene constante dependiente de $v$. Los radios y mayorantes analíticos
sí requieren cotas. \emph{Grados intermedios}, partes 04 y 05, construye radios
$c/v^8$ para esta comparación y obtiene un coste $(Cv)^{A_1v}$ por soporte.
El enlace correspondiente a pesos agrupados exige comparar ese coste con el
presupuesto de su norma y la ley de llegada. El hecho de que el cofactor coincida
en cada cara no da por sí solo sumabilidad al juntar todos los soportes.
